\documentclass[11pt]{report}
\input{../../../report.tex}
\title{Algebra Exercises}
\usepackage{titlepageBU}
\usepackage{xparse}
\courseID{CAS MA 741}
\college{College of Arts and Sciences}
\department{Department of Mathematics and Statistics}
\date{Winter Break, 2024/2025}
\usepackage{tikz-cd}
\DeclareMathOperator{\Stab}{Stab}
\DeclareMathOperator{\rank}{rk}
\declaretheoremstyle[headfont=\bfseries, bodyfont=\normalfont]{x}
\declaretheorem[numberwithin=section, style=x, name=Exercise]{exercise}
\declaretheorem[numbered=no, style=proof, name=Solution]{solution}
\begin{document}
\makereport
\stepcounter{chapter}
\stepcounter{chapter}
\stepcounter{chapter}
\stepcounter{chapter}
\stepcounter{chapter}
\stepcounter{chapter}
\stepcounter{chapter}
\chapter{Linear Algebra, Reprise}
\section{Functors}
\begin{exercise}
	Let $\mathscr{F} : \mathsf{C}  \to \mathsf{D} $ be a covariant functor, and let $\mathsf{C} $ and $\mathsf{D} $ have products. Show there exists a unique morphism $\mathscr{F} (A\times B)\to \mathscr{F} (A)\times \mathscr{F} (B)$ making the relevant diagram commute.

	If $\mathsf{D} $ has coproducts (denoted $A \amalg B$) and $\mathscr{G} : \mathsf{C}  \to \mathsf{D} $ is contravariant, show there exists a unique morphism $\mathscr{G}(A)\amalg \mathscr{G} (B)\to \mathscr{G} (A\times B)$.
\end{exercise}
\begin{solution}
	By definition of a functor, $\mathscr{F} (A\times B)$ is an object in $\mathsf{D} $, and
	\[\begin{tikzcd}[row sep = 2.5em, column sep = 0em]
		&\mathscr{F} (A\times B)\arrow[dr, "\mathscr{F} (\pi_B)"]\arrow[dl, "\mathscr{F} (\pi_A)"']\\
		\mathscr{F} (A)&&\mathscr{F} (B)
	\end{tikzcd}\]
	is a diagram in $\mathsf{D} $. Therefore, by universal property of products, the diagram
	\[\begin{tikzcd}[row sep = 2.5em, column sep = 2.5em]
		&&\mathscr{F} \left( A \right) \\
		\mathscr{F} (A\times B)\arrow[urr, bend left, "\mathscr{F} (\pi_A)"]\arrow[drr, bend right, "\mathscr{F} (\pi _B)"]\arrow[r, dashed]&\mathscr{F} (A)\times \mathscr{F} (B)\arrow[ur, "\pi_{\mathscr{F} (A)} "]\arrow[dr,"\pi_{\mathscr{F} (B)} "]\\
																		    &&\mathscr{F} (B)
	\end{tikzcd}\]
	commutes.

	Now suppose $\mathscr{G} : \mathsf{C}  \to \mathsf{D} $ is contravariant and $\mathsf{D} $ has coproducts. Since there exist projections $\pi_A: A\times B \to A$, $\pi_B : A\times B \to B$ in $\mathsf{C} $, and since $\mathscr{G} $ is contravariant, there exist morphisms
	\[
		\mathscr{G} (\pi_A): \mathscr{G} (A)\to \mathscr{G} (A\times B),\qquad \mathscr{G} (\pi_B): \mathscr{G} (B) \to \mathscr{G} (A\times B)
	\]
	in $\mathsf{D} $. Since $\mathsf{D} $ has coproducts, by the universal property of coproducts,
	\[\begin{tikzcd}[row sep = 2.5em, column sep = 2.5em]
		\mathscr{G} (A)\arrow[dr, "i_A"]\arrow[drr, bend left, "\mathscr{G} (\pi_A)"]\\
		&\mathscr{G} (A)\amalg \mathscr{G} (B)\arrow[r, dashed]&\mathscr{G} (A\times B)\\
		\mathscr{G} (B)\arrow[ur, "i_B"]\arrow[urr, bend right, "\mathscr{G} (\pi_B)"]
	\end{tikzcd}\]
	exists in $\mathsf{D} $ and commutes.
\end{solution}
\begin{exercise}
	Let $\mathscr{F} : \mathsf{C}  \to \mathsf{D} $ be a fully faithful functor. If $A,B$ are in $\mathsf{C} $, show that $A\cong B$ if and only if $\mathscr{F} (A)\cong \mathscr{F} (B)$.
\end{exercise}
\begin{solution}
	trivial, do later
\end{solution}
% 1.3 trivial too, 1.4 idk what its asking yet

\begin{exercise}
	Formalize the notion of a presheaf of abelian groups on a topological space $T$. If $\mathscr{F} $ is a presheaf on $T$, elements of $\mathscr{F} (U)$ are called \emph{sections} of $\mathscr{F} $ on $U$. The homomorphism $\rho_{UV} : \mathscr{F} (U) \to \mathscr{F} (V)$ induced by an inclusion $V\subseteq U$ is called the \emph{restriction map}.

	Show that a presheaf is obtained by letting $\mathscr{C} (U)$ be the additive abelian group of complex valued functions on $U$, with restriction of sections defined by ordinary restriction of functions.

	For this presheaf, prove that one can uniquely glue sections agreeing on overlapping sets. That is, if $U,V$ are open sets and $s_U\in \mathscr{C} (U)$, $s_V\in \mathscr{C} \left( V \right) $ agree after restriction to $U \cap V$, prove that there exists a unique $s\in \mathscr{C} (U \cup V)$ such that $s$ restricts to $s_U$ on $U$ and $s_V$ on $V$.
\end{exercise}
\begin{solution}
	Formalizing the notion is very simple. I will change notation slighty: let $(X,\mathcal{T} )$ be a topological space, and define the category $\mathsf{T} $ such that $\Obj(\mathsf{T} ) =\mathcal{T} $ and
	\[
		\Hom_{\mathsf{T} }(U, V) =
		\begin{cases}
			\varnothing ,&U\not\subseteq V\\
			\left\{ * \right\} ,&U\subseteq V
		\end{cases}
	.\]
	Then a presheaf $\mathscr{F} $ on $X$ is a contravariant functor $\mathscr{F} : \mathsf{T}  \to \mathsf{Ab} $.

	Let $\mathscr{C} : \mathsf{T}  \to \mathsf{Ab} $ be the contravariant assignment given by
	\[
		\mathscr{C} (U)=(C(U,\mathbb{C} ),+)
	.\]
	That is, $U$ maps to the set of continuous functions $f : U \to \mathbb{C} $, viewed as a group under addition. Morphisms $i : U \hookrightarrow V$ are defined as
	\[
		\mathscr{C} (i)(f)=\rho_{VU} (f)=f|_{U}
	.\]
	That is, the restriction to $U$ of a section $f$ of $\mathscr{C} $ over $V$ is the standard function restriction. Verification that the group given is abelian is immediate, so now we show that the restriction morphisms are group homomorphisms. Note that if $f,g\in \mathscr{C} (V)$ and $U\subseteq V$, then
	\[
		\rho_{VU} (f+g)(x)=(f+g)|_{U}(x)\coloneqq f|_U(x)+g|_U(x)=\rho_{VU} (f)(x)+\rho_{VU} (g)(x)
	.\]
	Therefore, this is a group homomorphism by definition of the operation. Next, we show identity maps map to identity maps. Note that for any $f : U \to U$,
	\[
		\rho_{UU} (f)=f_U=f	
	.\]
	Therefore, $\rho_{UU} =1_{C(U)} $. Finally, we wish to show that if $U\subseteq V\subseteq W$, then $\rho_{WU} =\rho_{WV} \rho_{VU} $. Let $f : W \to \mathbb{C} $. We have
	\[
		\rho_{WU} (f)=f_{U} ,\qquad \rho_{WV}\rho_{VU} (f)=\rho_{VU}f_V=f_U
	.\]
	Therefore, $\mathscr{C} $ is a contravariant functor, and thus a presheaf of abelian groups on $X$.

	Finally, the gluing condition. Let $s_U\in \mathscr{C} (U)$ and $s_V\in \mathscr{C} (V)$ agree on $U \cap V$. We can define a function $s\in \mathscr{C} (U \cup V)$ as the map
	\[
		s(x)=
		\begin{cases}
			s_U(x),&x\in U\\
			s_V(x),&x\in V
		\end{cases}
		
	.\]
	Note that this is well-defined because if $x\in U$ and $x\in V$, then $s_U(x)=s_V(x)$. Moreover, it is manifestly unique since it is uniquely determined by $s_U$ and $s_V$, and since they are continuous, $s$ is clearly continuous as well due to the fact that $s_U$ and $s_V$ agree on $U \cap V$. Therefore, $\mathscr{C} $ is a sheaf of abelian groups on $X$.
\end{solution}	


% 1.14
\begin{exercise}
	Verify that in $\Rmod{R} $, the equalizer of $\phi : A \to B$ and $0: A \to B$ is the kernel $\ker \phi $.
\end{exercise}
\begin{solution}
	Verification is simple. Note that for all objects $X$, the diagram
	\[\begin{tikzcd}[row sep = 3.5em, column sep = 2.5em]
		A\arrow[rr, bend left, "\phi "]\arrow[rr, bend right, "0"]&&B\\
									  &\varprojlim \mathscr{F} \arrow[ul, "f_A"]\arrow[ur, "f_B"']\\
									  &X\arrow[luu, bend left, "x_A"]\arrow[u, dashed]\arrow[ruu, bend right, "x_B"']
	\end{tikzcd}\]
	commutes. Moreover, $f_B=0f_A=\phi f_A$. However, since $0f_A=0$, we have $\phi f_A=0$, and $f_B=0$. We also have $x_B=0x_A=\phi x_A$, and thus $x_B=0$. We can rewrite this diagram as
	\[\begin{tikzcd}[row sep = 2.5em, column sep = 2.5em]
		X\arrow[r, "x_a"]\arrow[dr, dashed]\arrow[rr, bend left, "0"]&A\arrow[r, "\phi "]&B\\
									     &\varprojlim \mathscr{F} \arrow[ur, "0"]\arrow[u, "f_A"]
	\end{tikzcd}\]
	Note that this is precisely the universal property satisfied by the kernel, since by commutativity of the diagram, we have that $\phi f_A=0$, and we can actually drop the 0 morphism $\varprojlim \mathscr{F} \to B$ while retaining precisely the same information. Therefore, the same proof given that kernels satisfy this universal property in $\Rmod{R} $ suffices here to show that $\varprojlim \mathscr{F} $ exists, and is the kernel.
\end{solution}

% 1.15
\begin{exercise}
	Let $\mathsf{I} $ be the category consisting of elements of $\mathbb{Z}^+$ as objects, and a unique morphism $i\to j$ if and only if $i\leq j$. Let $\mathscr{F} : \mathsf{I}  \to \Rmod{R} $ be a functor, consisting of a choice of object $A_i$ for each $i$, and a morphism $\phi_{ji} : A_i \to A_j$ for all $i,j$. Note that an element of
	\[
		\prod_{i} A_i
	\]
	is a sequence $(a_i)_{i>0} $. Call a sequence $(a_i)$ coherent if $\phi_{i+1,i} (a_{i+1} )=a_i$ for all $i$. Let $A$ be the submodule of $\prod_{i} A_i$ consisting of all coherent sequences. Let $\phi_i : A \to A_i$ be the restriction of the canonical projection $\pi_i : \prod_{i} A_i \to A_i$ to $A$. Show that $A$, together with the morphisms $\phi_i$, satisfy
	\[
		\varprojlim_i A_i=A
	.\]
\end{exercise}
\begin{solution}
	First, I want to show $A$ is truly an $R$-module. Let $(a_i),(b_i)\in A$. Then
	\[
		(a_i)-(b_i)=(a_i-b_i)
	.\]
	Note that (with the mild abuse of notation), we have
	\[
		\phi_{i+1,i} (a_{i+1} -b_{i+1}) =\phi_{i+1,i}(a_{i+1})-\phi_{i+1,i}(b_{i+1} )=(a_i)-(b_i)=(a_i-b_i)
	.\]
	Additionally, for $r\in R$,
	\[
		r(a_i)=(ra_i)
	.\]
	Thus,
	\[
		\phi_{i+1,i} (ra_{i+1} )=r\phi_{i+1,i} (a_{i+1} )=r(a_i)=(ra_i)
	.\]
	Therefore, $A$ is an $R$-module, as perscribed. Now we show the relevant diagram commutes. Let $n\in\mathbb{Z} ^+$. We have $\phi_n(a_i)=a_n$. We also have
	\[
		\left( \phi_{n+1,n}\phi_{n+1} \right) (a_i)=\phi_{n+1,n} (a_{n+1} )=a_n
	.\]
	Therefore, the diagram
\[\begin{tikzcd}[row sep = 2.5em, column sep = 2.5em]
	&&A\arrow[dll, "\ldots "']\arrow[dl, "\phi_4"]\arrow[d, "\phi_3"]\arrow[dr, "\phi _2"]\arrow[drr, "\phi _1"]\\
	\cdots \arrow[r, "\phi_{45} "]&A_4\arrow[r, "\phi_{34} "]&A_3\arrow[r, "\phi_{23} "]&A_2\arrow[r, "\phi_{12} "]&A_1
\end{tikzcd}\]
clearly commutes. Finally, $A$ is manifestly final due to being the subset of a product and inheriting the canonical projections.
\end{solution}

\end{document}
